\documentclass[10pt,a4paper]{amsart}
\usepackage[margin=19mm]{geometry}
\usepackage{amsmath,amssymb,mathtools}
\usepackage[scr=rsfs,cal=euler]{mathalfa}
\usepackage{xfrac}
\usepackage[unicode,hidelinks,bookmarks=false]{hyperref}
\newcommand{\ie}{i.e.\ }
\newcommand{\cf}{cf.\ }
\newcommand{\resp}{respectively\ }
\newcommand{\Subsection}{\subsection}
\providecommand{\nobreakdash}{\nobreak\hbox{-}}
\setlength{\emergencystretch}{2em}
\multlinegap0pt
\usepackage{amssymb}
\usepackage{float}
\newtheorem{theorem}{Theorem}
\title{Symmetric Sudoku-Type Games from Perfect Codes}
\author{Junmin An, Jae-Hyun Baek, Keon-Hwi Kim, Haeun Lim, Jon-Lark Kim}
\date{Source: arXiv:2605.09617v1 (CC BY 4.0)}
\begin{document}
\maketitle

\noindent\emph{Source and licence.} Symmetric Sudoku-Type Games from Perfect Codes. Version arXiv:2605.09617v1 (CC BY 4.0), distributed by arXiv under the Creative Commons Attribution 4.0 International licence, http://creativecommons.org/licenses/by/4.0/, as recorded in the arXiv licence metadata for this exact version and shown on the abstract page https://arxiv.org/abs/2605.09617v1. Full licence notice: \texttt{licenses/arxiv-2605.09617-CC-BY-4.0.txt}.\par\smallskip

\noindent\textit{Editorial scope.} Counted unit: the article's theorem theorem (source0.tex lines 407--420). The article's complete original proof of it is delivered with it (source0.tex lines 421--481, one \texttt{proof} environment of 61 lines). No mathematics is written, repaired or re-derived: the statement and the proof are the article's own lines, in the article's own notation, and no retained line is substituted or re-flowed.  No further source-local context is needed: every cross-reference of every retained line resolves inside the two delivered spans.  Nothing was omitted from any retained span, so the package carries no omission record.  No retained line contains a citation, so the wrapper carries no bibliography and no bibliography entry.  No retained line contains a figure environment, an image-inclusion command or drawing code, so no image is delivered and the package declares no asset dependency.  Editorial formatting only, in addition: an AMS \texttt{amsart} preamble that restates the article's own declarations and macros used by the retained lines, namely the environments \texttt{theorem} on separate counters, together with the standard packages those lines need.  The retained environments print in this document as Theorem 1 in order of appearance, whereas the article prints the same items with the numbering of its own full document.  The complete native source of the article as distributed in the arXiv e-print for this exact version is bundled unchanged as \texttt{sources/200241/source0.tex}.
\begin{theorem}
Let $\mathcal{C}_0$ be a linear $t$-error-correcting diameter perfect code with generator matrix $G_0$ and $\mathcal{C}=\mathbf{x}+\mathcal{C}_0\subseteq\mathbb{Z}_n^2$ for some $\mathbf{x}=(x_1, x_2)\in\mathbb{Z}_n^2$. Suppose $\mathcal{C}$ is either of Case I or Case II.
	\begin{enumerate}
	\item[(i)] If $\mathcal{C}$ is of Case I, let
	\[
	\mathcal{G}_\mathcal{S}=\langle \tau_1^{t+1}\tau_2^{t+1},\, \tau_2^{2(t+1)},\,\tau_2^{2x_2+1}s,\, \tau_1^{2x_1+2}\tau_2^{2x_2+1}r^2\rangle.
	\]
	\item[(ii)] If $\mathcal{C}$ is of Case II and $G_0=[a~~b]$, let
	\[
	\mathcal{G}_\mathcal{S}=\langle \tau_1^{a}\tau_2^{b},\, \tau_1^{2x_1+2}\tau_2^{2x_2+1}r^2\rangle.
	\]
	\end{enumerate}
	Then, for every $S\in \mathcal{S}$ and $g\in \mathcal{G}_\mathcal{S}$, we have $g\cdot S\in \mathcal{S}$.
\end{theorem}

\begin{proof}
We show that, in each case, the action of every generator of $\mathcal{G}_\mathcal{S}$ on a Sudoku grid $S \in \mathcal{S}$ preserves orthogonality with the corresponding palette grid $\mathcal{I}_{\mathcal{C}, \mathcal{A}}$.
\begin{enumerate}
\item[(i)] Let us assume that $\mathcal{C}$ is of Case I. We consider each generator of $\mathcal{G}_\mathcal{S}$ separately. Take any $\mathbf{c}=(c_1, c_2)\in\mathcal{C}$ and $S\in\mathcal{S}$.
	\begin{itemize}
	\item[(1)] We have
		\[
		(\tau_1^{t+1}\tau_2^{t+1}(S))_{c_1, c_2}=S_{c_1-(t+1), c_2-(t+1)}
		\]
		and
		\[
		(\tau_2^{2(t+1)}(S))_{c_1, c_2}=S_{c_1, c_2-2(t+1)}.
		\]
		Since both $(c_1-(t+1), c_2-(t+1))$ and $(c_1, c_2-2(t+1))$ are codewords of $\mathcal{C}$, $\tau_1^{t+1}\tau_2^{t+1}(S)$ and $\tau_2^{2(t+1)}(S)$ are orthogonal to $\mathcal{I}_{\mathcal{C}, \mathcal{A}}$.
	\item[(2)] Note that
		\[
		(\tau_2^{2x_2+1}(s(S)))_{c_1, c_2}=(s(S))_{c_1, c_2-(2x_2+1)}=S_{c_1, n-c_2+2x_2}.
		\]
		Since $\mathbf{c}=\mathbf{x}+\alpha(t+1, t+1)+\beta(0, 2(t+1))$ for some $\alpha, \beta\in\mathbb{Z}_n$, it follows that
		\[
		(c_1, n-c_2+2x_2)=\mathbf{x}+\alpha(t+1, t+1)-(\alpha+\beta)(0, 2(t+1))
		\]
		is also a codeword of $\mathcal{C}$. So $(\tau_2^{2x_2+1}s(S))$ is orthogonal to $\mathcal{I}_{\mathcal{C}, \mathcal{A}}$.
	\item[(3)] Since a $180^\circ$ rotation reverses the positions of the two entries in the core of each translated anticode on the grid, it suffices to show that
		\[
		(\tau_1^{2x_1+2}\tau_2^{2x_2+1} r^2(S))_{c_1, c_2} = S_{c_1'+1, c_2'}
		\]
		and
		\[
		(\tau_1^{2x_1+2}\tau_2^{2x_2+1} r^2(S))_{c_1+1, c_2} = S_{c_1', c_2'}
		\]
		for some codeword $\mathbf{c}' = (c_1', c_2') \in \mathcal{C}$. Let $x'=2x_1+2$ and $x''=2x_2+1$. Then we compute:
		\[
			(\tau_1^{x'}\tau_2^{x''} (r^2(S)))_{c_1, c_2}=(r^2(S))_{c_1-x', c_2-x''}=S_{2x_1-c_1+1, 2x_2-c_2}.
		\]
		If we take $\mathbf{c}'=2\mathbf{x}-\mathbf{c}\in\mathcal{C}$, then we have
		\[
		(\tau_1^{2x_1+2}\tau_2^{2x_2+1} r^2(S))_{c_1, c_2} = S_{c_1'+1, c_2'}.
		\]
		Similarly,
		\[
			(\tau_1^{x'}\tau_2^{x''} r^2(S))_{c_1+1, c_2}=S_{2x_1-c_1, 2x_2-c_2}=S_{c_1', c_2'}.
		\]
		Therefore, $\tau_1^{2x_1+2-n} \tau_2^{2x_2+1-n} r^2(S)$ is orthogonal to $\mathcal{I}_{\mathcal{C}, \mathcal{A}}$.
	\end{itemize}
	Hence $g \cdot S \in \mathcal{S}$ for every generator $g \in \mathcal{G}_\mathcal{S}$, and the group $\mathcal{G}_\mathcal{S}$ acts on $\mathcal{S}$ with respect to the code $\mathcal{C}$.
\item[(ii)] Suppose that $\mathcal{C}$ is of Case II. As in the previous case, we verify that each generator of $\mathcal{G}_\mathcal{S}$ maps any $S \in \mathcal{S}$ to another element in $\mathcal{S}$. Take any $\mathbf{c} = (c_1, c_2) \in \mathcal{C}$ and $S \in \mathcal{S}$.

Note that
\[
(\tau_1^{a}\tau_2^{b}(S))_{c_1, c_2}=S_{c_1-a, c_2-b}.
\]
Since $\mathbf{c}-(a, b)\in\mathcal{C}$, it follows that $\tau_1^{a}\tau_2^{b}(S)$ is orthogonal to $\mathcal{I}_{\mathcal{C}, \mathcal{A}}$.

Note that the proof of (i)-(3) does not depend on whether $\mathcal{C}$ is of Case I or Case II. Therefore, as shown in (i), we have
\[
\mathcal{I}_\mathcal{C}\perp\tau_1^{2x_1+2} \tau_2^{2x_2+1} r^2(S).
\]
Thus $g\cdot S\in\mathcal{S}$ for all $g\in \mathcal{G}_\mathcal{S}$. This completes the proof.
\end{enumerate}
\end{proof}

\end{document}
